% !TEX root = ../../mixed-sklyanin-cubics.tex
% Current result: 2026-11-22 / v51.
\section{対合を保つ三対象生成子の Hom 指標}
\label{mix:sec:equivariant-window}

定理\ref{cmo:thm:exact-sixty-four}の指定係数に対して
\[
 B=\mathcal B_{(1,4,9)}\otimes_{\C[u]}\mathcal B_{(16,25,49)},
 \qquad w=\Omega_1+100\Omega_2,\qquad A_0=B/(w)
\]
と置く。六つの active 生成元を固定し、$u\mapsto-u$ とする対合を
$\sigma$ と書く。二次関係と $w$ は $u$ の偶数冪だけを含むので、
$\sigma$ は $(B,w)$ の自己同型である。
本節では、$\operatorname{Ku}_{\mathrm{coh}}(A_0)$ の有限生成子の
自己 Hom における対合の作用を計算する。
これは Hochschild homology における対合の作用とは別の計算である。

内部 shift は $M(j)_q=M_{q+j}$ とし、指定された三次周期を
$(3)\simeq[2]$ とする。augmentation 加群 $\C$ に自明な
$\sigma$ 作用を入れ、
\[
 K=\operatorname{stab}(\C),\qquad
 G=K\oplus K(1)\oplus K(2),\qquad
 \Gamma=R\operatorname{End}(G)
\]
にそれから誘導される同変構造を入れる。
各内部 shift の同変構造に、追加の符号指標は掛けない。
$B$ は大域次元七の Koszul AS-regular 代数であり、$w$ は正則中心三次元である。
$A_0$ は左右 noetherian な AS-Gorenstein 代数で、injective dimension は六、
Gorenstein parameter は $a=4$ である。
命題\ref{cmo:prop:proper-contraction}の証明で確認した
$\chi$ 条件、有限 cohomological dimension、線束 Hom の公式と七本の線束の生成性を使う。
ここで適用するのは三対象 window 定理\paperref{main}{tw:thm:window}の
証明中の体上の section--window 部分であり、同節の可換特殊化の仮定ではない。
実際、その証明で用いた非可換 Gorenstein 代数の有界 Orlov 分解は
上の入力だけで適用でき、
\[
 \Psi=\operatorname{stab}\circ(R\omega(-))_{\ge-3}
\]
は四つの例外線束を除いた線束残余圏上で完全忠実である。
これは体上の Orlov 商そのものによる完全忠実性であり、
可換ファイバーとの Hom 比較または導来 Nakayama の適用を要しない。

生成性も同じ証明の有限 filtration で確認できる。
$4\le j\le6$ に対して $T_j=A_0(j)_{<-3}$ とすると、
\[
 \Psi(\operatorname{pr}\mathcal O(j))\simeq\operatorname{stab}(T_j)[-1],
 \qquad
 \operatorname{gr}T_j=\bigoplus_{r=0}^{j-4}(A_0)_r\otimes\C(j-r).
\]
従って三角的に $K(4),K(5),K(6)$ が生成され、$(3)\simeq[2]$ により
その thick closure は $G$ のものと一致する。
線束側の properness は既存の Hom 公式から従い、
$G$ の Hom の有限性は以下でも直接計算する。
定理\ref{cmo:thm:exact-sixty-four}により $A_0$ は smooth なので、
同圏は全 coherent 残余圏と一致する。
section、内部切断、augmentation stabilization の各操作は
$\sigma$ と可換であり、この比較は同変に選べる。

\begin{proposition}[三対象生成子の stable Hom の対合指標]
\label{cmo:prop:equivariant-window-hom}
$W_i=\operatorname{Tor}^B_i(\C,\C)$ とし、$W_i$ の内部次数を $i$ とする。
$0\le i\le7$ に対して
\begin{equation}\label{cmo:eq:koszul-involution-character}
 \dim W_i=\binom7i,\qquad
 \sum_{i=0}^7\operatorname{Tr}(\sigma\mid W_i)z^i
 =(1+z)^6(1-z).
\end{equation}
$0\le a,b\le2$、$d=b-a$ とすると、$\sigma$ 表現として
\begin{equation}\label{cmo:eq:stable-window-character}
 H^rR\operatorname{Hom}(K(a),K(b))
 \cong
 \bigoplus_{\substack{0\le i\le7,\ p\in\mathbb Z\\
                       i+3p=-d,\ r=i+2p}}W_i^*.
\end{equation}
この同型は加群の直和分解と対合の指標を述べるものであり、
Yoneda 積または高次 $A_\infty$ 積の表示を主張しない。

次の表が全ての非零 Hom を与える。
$h_+$ と $h_-$ は、それぞれ $+1$ と $-1$ 固有空間の次元である。
\begin{equation}\label{cmo:eq:window-hom-table}
\begin{array}{c|rrr|rrrr}
 d&r&i&p&\dim&\operatorname{Tr}\sigma&h_+&h_-\\\hline
 -2& 2&2& 0&21& 9&15& 6\\
 -2& 3&5&-1&21&-9& 6&15\\\hline
 -1& 1&1& 0& 7& 5& 6& 1\\
 -1& 2&4&-1&35&-5&15&20\\
 -1& 3&7&-2& 1&-1& 0& 1\\\hline
  0& 0&0& 0& 1& 1& 1& 0\\
  0& 1&3&-1&35& 5&20&15\\
  0& 2&6&-2& 7&-5& 1& 6\\\hline
  1& 0&2&-1&21& 9&15& 6\\
  1& 1&5&-2&21&-9& 6&15\\\hline
  2&-1&1&-1& 7& 5& 6& 1\\
  2& 0&4&-2&35&-5&15&20\\
  2& 1&7&-3& 1&-1& 0& 1
\end{array}
\end{equation}
特に $H^*\Gamma$ は合計 $384$ 次元であり、
\begin{equation}\label{cmo:eq:window-total-table}
\begin{array}{c|rrr}
 r&\dim H^r\Gamma&\dim(H^r\Gamma)^+&\dim(H^r\Gamma)^-\\\hline
 -1& 7& 6& 1\\
  0&80&48&32\\
  1&162&84&78\\
  2&112&48&64\\
  3&23&6&17
\end{array}
\end{equation}
となる。
\end{proposition}

\begin{proof}
まず $u$ は $B$ 上正則で、$B/(u)=S(1,-1,-1)^{\otimes2}$ の上では
$\sigma$ が恒等である。この商の Hilbert 級数は $(1-z)^{-6}$ である。
$B(-1)\xrightarrow{u}B$ を同変にするために源へ符号指標を掛けると、
次数別指標の完全列から
\[
 H_B^\sigma(z)=\frac1{(1-z)^6(1+z)},\qquad H_B(z)=\frac1{(1-z)^7}
\]
を得る。標準 Koszul 分解 $P_i=B(-i)\otimes W_i$ は同変であり、
その Euler 恒等式は
\[
 H_B^\sigma(z)\sum_{i=0}^7(-1)^i\operatorname{Tr}(\sigma\mid W_i)z^i=1
\]
である。通常の Hilbert 級数にも同じ恒等式を適用すれば
式\eqref{cmo:eq:koszul-involution-character}を得る。

次に、この指標を実際の stable Ext へ移す。
補題\paperref{main}{mbnp:lem:shamash}の higher homotopies を
同変に構成する。$s_0=d$ とし、homological degree を下添字で表すと
\[
 s_q:P_i\longrightarrow P_{i+2q-1}(3q),\qquad
 \sum_{a+b=q}s_as_b=
 \begin{cases}
  w\operatorname{id},&q=1,\\
  0,&q\ne1
 \end{cases}
\]
である。$q=1$ では $w$ の augmentation 上の作用が零なので
nullhomotopy が存在し、その $\sigma$ による共役との平均を取る。
$s_1,\ldots,s_{q-1}$ を同変に選んだ後では、次の式の右辺
\[
 d s_q+s_qd=-\sum_{a=1}^{q-1}s_as_{q-a}
\]
は同変である。右辺が閉じていることと、
$\operatorname{End}_B(P_\bullet)$ の負の cohomology の消滅により解が存在する。
その解だけを共役と平均すれば、この線型方程式を保ったまま
$s_q$ を同変にできる。
既に選んだ低次の写像は変更しないので、以前の二次恒等式も保たれる。
これは各 $s_q$ を独立に平均する議論ではない。
分解の長さが七なので $q>4$ の写像は零である。

$P_i$ と $P_{i+2q-1}(3q)$ の生成次数から、$s_q$ の行列係数は
次数 $q+1$ を持つ。従って $q\ge1$ の係数は正次数であり、
Koszul 微分の係数も次数一である。
同補題で構成された $A_0$ 上の分解
\begin{equation}\label{cmo:eq:equivariant-minimal-shamash}
 F_n=\bigoplus_{q\ge0}
       A_0(-i-3q)\otimes W_i,\qquad i=n-2q,\quad 0\le i\le7
\end{equation}
は同変な minimal 分解となる。
完全性には同補題の filtration を用いる。
その第一頁は $\operatorname{Tor}^B(\C,A_0)
=\C\oplus\C(-3)[1]$ の各周期コピーであり、
$s_1$ が隣接する二項を同型で相殺するため、augmentation だけが残る。
従って、単なる Hilbert 級数の一致から分解の完全性を推論してはいない。

十分高い次数で $F_{n+2}=F_n(-3)$ となる。
この同定は $w$ による指定周期から来ており、$\sigma(w)=w$ と
同変 higher homotopies の構成から、その周期変数 $t$ の指標は $+1$ である。
周期変数は cohomological degree 二、分解の生成元の内部次数三を持つ。
従って Ext の標的内部 shift の表記では、その双次数は $(2,-3)$ である。
$A_0$ は正則中心商で Iwanaga--Gorenstein であるから、
この周期尾部は augmentation の complete resolution を与える。
stabilization と graded 特異圏の比較
（補題\paperref{main}{tw:lem:stabilization}の体上の代数的余核比較）により、
周期を反転した分解で $K$ の stable Hom を計算できる。
minimality によって $\operatorname{Hom}_{A_0}(-,\C)$ の微分は零となり、
得られる双次数付き表現は
\[
 \bigoplus_{i=0}^7 W_i^*\otimes\C[t,t^{-1}]
\]
である。
ここでもこの式を積付き代数の同型としては用いない。

$K(a)$ から $K(b)$ への次数零の内部 Hom を取る条件は
$i+3p=-(b-a)$ であり、cohomological degree は $r=i+2p$ である。
これで式\eqref{cmo:eq:stable-window-character}を得る。
対合の双対表現の指標は元と等しい。
式\eqref{cmo:eq:koszul-involution-character}の係数は
\[
 (1,5,9,5,-5,-9,-5,-1)
\]
なので、$h_\pm=(\binom7i\pm\operatorname{Tr}(\sigma\mid W_i))/2$ を代入すれば
表\eqref{cmo:eq:window-hom-table}となる。
各 $d$ は三つの対象の間に $3-|d|$ 回現れるため、
その重複度で集計すると表\eqref{cmo:eq:window-total-table}を得る。
\end{proof}

\begin{corollary}[通常および対合付き Hom Euler 行列]
\label{cmo:cor:equivariant-window-euler}
順序 $(K,K(1),K(2))$ に関して
\[
 \begin{aligned}
 E_{ab}&=\sum_r(-1)^r\dim H^rR\operatorname{Hom}(K(a),K(b)),\\
 E^\sigma_{ab}&=\sum_r(-1)^r\operatorname{Tr}
       \bigl(\sigma\mid H^rR\operatorname{Hom}(K(a),K(b))\bigr)
 \end{aligned}
\]
と置くと、
\begin{equation}\label{cmo:eq:equivariant-window-euler}
 E=\begin{pmatrix}-27&0&27\\27&-27&0\\0&27&-27\end{pmatrix},
 \qquad
 E^\sigma=\begin{pmatrix}-9&18&-9\\-9&-9&18\\18&-9&-9\end{pmatrix}.
\end{equation}
不変 Hom の Euler 行列は
\[
 \frac{E+E^\sigma}{2}
 =\begin{pmatrix}-18&9&9\\9&-18&9\\9&9&-18\end{pmatrix}
\]
である。
\end{corollary}
\begin{proof}
表\eqref{cmo:eq:window-hom-table}の各行を交代和する。
既存の線束計算との整合は、独立に次のように確認できる。
三対象 window の三角的 filtration
（式\papereqref{main}{tw:eq:triangular}）を $a=4$ に適用すると、
系\ref{cmo:cor:rational-euler-comparison}の類は
\[
 e_4=-[K(1)],\qquad e_5=-7[K(1)]-[K(2)]
\]
となる。従って $(e_4,e_5-7e_4)$ の Euler 行列は
式\eqref{cmo:eq:equivariant-window-euler}の隣接二対象の小行列
$27\bigl(\begin{smallmatrix}-1&0\\1&-1\end{smallmatrix}\bigr)$ に一致する。
\end{proof}

\begin{remark}[有限な Hom 計算と Hochschild 計算の境界]
\label{cmo:rem:window-not-hochschild}
命題\ref{cmo:prop:equivariant-window-hom}は、
対合を保つ minimal $A_\infty$ 模型の underlying cohomology を
$384$ 次元に確定する。しかし、その高次積と Hochschild 微分はまだ計算していない。
特に $E^\sigma$ は対象間の Hom の supertrace であり、
$HH_*(\operatorname{Ku}_{\mathrm{coh}}(A_0))$ の supertrace ではない。

この区別は bar 複体を有限長で切る際にも必要である。
$H^1R\operatorname{Hom}(K,K)$ は三十五次元で、単位を含まない。
従って正規化 Hochschild bar 複体でこれらを suspension すると次数零となり、
同じ総次数に任意の長さの composable な bar 語が存在する。
内部 shift も全て零なので、三次周期の grading からこの長さは制限されない。
有限個の $m_r$ または有限 bar 長の計算だけでは、
残りを縮約する証明なしに Hochschild homology の計算にならない。
必要なのは対角の有限解消、正当化された bar 縮約、または別の同変比較である。
非可換側の $HH_{\pm1}$ の固有次元が可換候補の $(5,16)$ と一致するかは、
本節の計算からは結論しない。

\texttt{computations/verify\_mixed\_window\_involution.py} は
指標多項式、全 Hom 表、集約表、Euler 行列を整数で再計算する。
さらに線束の Hilbert 多項式と例外射影から Euler 行列を独立に作り、
上の window の基底変換と照合する。
この検算に Hochschild 指標の未証明の値は含めない。
\end{remark}
